\documentclass[10pt]{article}
\usepackage[T1]{fontenc}
\usepackage[utf8]{inputenc}
\usepackage[french]{babel}
\usepackage{amsmath,amssymb,mathrsfs,mathpazo}
\usepackage{geometry}
\geometry{a4paper,margin=2cm}
\usepackage{graphicx,array,tabularx,tabularray,colortbl,multirow,multicol,booktabs,enumitem,pifont,ulem,fancybox,diagbox,cancel}
\usepackage{pstricks,pstricks-add,pst-plot,pst-tree,pst-node,pst-text,pst-func,pst-math,pst-bspline,pst-3dplot,pst-eucl,pst-solides3d}
\let\pscalculate\undefined
\usepackage{pst-bezier}
\usepackage{tikz,pgfplots,tkz-tab}
\usetikzlibrary{arrows,arrows.meta,patterns,calc,positioning,trees,decorations.markings,shapes.geometric,shapes.misc,shapes.arrows,intersections,decorations.pathmorphing,decorations.pathreplacing,backgrounds,shadows,babel}
\pgfplotsset{compat=1.18}
\usepgfplotslibrary{fillbetween}
\usepackage[np]{numprint}
\usepackage{listings,makecell,pas-tableur,twoopt}
\let\barre\undefined
\usepackage[pstricks]{tablvar}
\definecolor{darkgreen}{RGB}{0,100,0}
\providecommand{\red}{\color{red}}
\providecommand{\blue}{\color{blue}}
\providecommand{\green}{\color{green}}
\providecommand{\black}{\color{black}}
\providecommand{\white}{\color{white}}
\providecommand{\gray}{\color{gray}}
\providecommand{\cyan}{\color{cyan}}
\providecommand{\magenta}{\color{magenta}}
\providecommand{\yellow}{\color{yellow}}
\providecommand{\violet}{\color{violet}}
\providecommand{\orange}{\color{orange}}
\providecommand{\euro}{\texteuro}
\providecommand{\EUR}{\texteuro}
\providecommand{\ieme}{\textsuperscript{e}}
\providecommand{\degres}{\ensuremath{^\circ}}
\providecommand{\degree}{\ensuremath{^\circ}}
\providecommand{\vv}[1]{\overrightarrow{#1}}
\providecommand{\Coord}[1]{#1}
\providecommand{\decosix}{}
\providecommand{\starredbullet}{\ensuremath{\bullet}}
\pagestyle{empty}
\setlength{\parindent}{0pt}
\usepackage{adjustbox}\DeclareUnicodeCharacter{2236}{:}\DeclareUnicodeCharacter{2190}{\ensuremath{\leftarrow}}
\providecommand{\R}{\mathbb{R}}
\providecommand{\N}{\mathbb{N}}
\providecommand{\Z}{\mathbb{Z}}
\providecommand{\Q}{\mathbb{Q}}
\providecommand{\D}{\mathbb{D}}
\providecommand{\C}{\mathbb{C}}
\providecommand{\e}{{\mathrm{e}}}
\providecommand{\vect}[1]{\overrightarrow{#1}}
\providecommand{\vectt}[1]{\overrightarrow{\mathrm{#1}}}
\providecommand{\vvt}[1]{\overrightarrow{\mathrm{#1}}}
\providecommand{\np}[2][]{{#2~#1}}
\providecommand{\vv}[1]{{\overrightarrow{#1}}}
\renewcommand{\d}{{\mathrm{d}}}
\providecommand{\degres}{{^\circ}}
\providecommand{\degre}{{^\circ}}
\providecommand{\blue}{\color{blue}}
\providecommand{\red}{\color{red}}
\providecommand{\enskip}{\quad}
\providecommand{\llbracket}{\lbrack\!\lbrack}
\providecommand{\rrbracket}{\rbrack\!\rbrack}
\providecommand{\euro}{{\text{€}}}
\providecommand{\no}{{\text{n°}}}
\providecommand{\interval}[2]{{[#1;#2]}}
\providecommand{\intervFO}[2]{{[#1;#2[}}
\providecommand{\dotfill}{\ldots}

\providecommand{\ds}{\displaystyle}
\providecommand{\textding}[1]{#1}
\providecommand{\Oij}{\(\mathrm{O};\vec{\imath},\vec{\jmath}\)}
\providecommand{\Oijk}{\(\mathrm{O};\vec{\imath},\vec{\jmath},\vec{k}\)}
\providecommand{\Ouv}{\(\mathrm{O};\vec{u},\vec{v}\)}
\providecommand{\intervFF}[2]{\ensuremath{[#1;#2]}}
\renewcommand{\euro}{\ifmmode\text{\texteuro}\else\texteuro\fi}
\DeclareUnicodeCharacter{00B0}{\ifmmode^{\circ}\else\textdegree\fi}
\DeclareUnicodeCharacter{20AC}{\ifmmode\text{\texteuro}\else\texteuro\fi}
% Énoncés : Éducation nationale. Transcription et corrigés : APMEP, Denis Vergès et contributeurs. Découpage et classement : CoopMaths, Sébastien Lozano, Éric Elter et contributeurs. Adaptation de la présentation : Maths Explorer.
% CC BY-SA : https://coopmaths.fr/www/about/
% https://coopmaths.fr/alea/static/bac/2023/tex/bac_2023_03_sujet1_etrangers_groupe1_2_cor.tex
\begin{document}
\noindent\textit{APMEP / CoopMaths. CC BY-SA. Présentation adaptée par Maths Explorer.}\par\medskip
\bigskip
\textbf{Partie A}

Interprétons l'énoncé :

si $n$ est un entier naturel, la phrase \og lorsque la trottinette est en bon état un lundi, la probabilité qu'elle  soit encore en bon état le lundi suivant est 0,9\fg{} correspond à : $P_{B_n}(B_{n+1}) = 0,9$.

la phrase suivante se traduit par : $P_{\overline{B_n}}(B_{n+1}) = 0,4$.

\begin{enumerate}
\item On a donc $p_1 = P(B_1) = P(B_0) \times P_{B_0}(B_{1}) = 1 \times 0,9 = 0,9$. On a $p_1 = 0,9$.

	Pour déterminer $p_2$, traçons un arbre :
	
	\hfill~
		
		\begin{tikzpicture}[xscale=1,yscale=1,baseline=(R)]
			
			\tikzstyle{fleche}=[->,>=latex,thick]
			\tikzstyle{noeud}=[fill=white,circle,inner sep=2pt]
			\tikzstyle{feuille}=[fill=white,circle,inner sep=2pt]
			\tikzstyle{etiquette}=[midway,fill=white, inner xsep=3pt, inner ysep=1.5pt]
			
			\def\DistanceInterNiveaux{3}
			\def\DistanceInterFeuilles{0.8}
			
			\def\NiveauA{(0)*\DistanceInterNiveaux}
			\def\NiveauB{(1)*\DistanceInterNiveaux}
			\def\NiveauC{(2)*\DistanceInterNiveaux}
			\def\InterFeuilles{(-1)*\DistanceInterFeuilles}
			
			\node[noeud] (R) at ({\NiveauA},{(1.5)*\InterFeuilles}) {$$};
			\node[noeud] (Ra) at ({\NiveauB},{(0.5)*\InterFeuilles}) {$B_1$};
			\node[feuille] (Raa) at ({\NiveauC},{(0)*\InterFeuilles}) {$B_2$};
			\node[feuille] (Rab) at ({\NiveauC},{(1)*\InterFeuilles}) {$\overline{B_2}$};
			\node[noeud] (Rb) at ({\NiveauB},{(2.5)*\InterFeuilles}) {$\overline{B_1}$};
			\node[feuille] (Rba) at ({\NiveauC},{(2)*\InterFeuilles}) {$B_2$};
			\node[feuille] (Rbb) at ({\NiveauC},{(3)*\InterFeuilles}) {$\overline{B_2}$};
			
			\draw[fleche] (R.east)--(Ra.west) node[etiquette] {$p_1 = 0,9$};
			\draw[fleche] (Ra.east)--(Raa.west) node[etiquette] {$0,9$};
			\draw[fleche] (Ra.east)--(Rab.west) node[etiquette] {$0,1$};
			\draw[fleche] (R.east)--(Rb.west) node[etiquette] {$0,1$};
			\draw[fleche] (Rb.east)--(Rba.west) node[etiquette] {$0,4$};
			\draw[fleche] (Rb.east)--(Rbb.west) node[etiquette] {$0,6$};
		\end{tikzpicture}
	\hfill~
	

Les évènements $B_1$ et $\overline{B_1}$ partitionnant l'univers, la loi des probabilités totales donne :

$p_2 = P(B_2) = P\Big(B_1 \cap B_2\Big) + P\Big(\overline{B_1} \cap B_2 \Big) = 0,9 \times 0,9 + 0,1 \times 0,4 = 0,81 + 0,04 = 0,85$.

Ce calcul confirme bien que l'on a : $p_2 = 0,85$.

	\item On va donc avoir un arbre très similaire au précédent : les probabilités conditionnelles restant les mêmes.

	\hfill~
			
			\begin{tikzpicture}[xscale=1,yscale=1,baseline=(R)]
				
				\tikzstyle{fleche}=[->,>=latex,thick]
				\tikzstyle{noeud}=[fill=white,circle,inner sep=2pt]
				\tikzstyle{feuille}=[fill=white,circle,inner sep=2pt]
				\tikzstyle{etiquette}=[midway,fill=white, inner xsep=3pt, inner ysep=1.5pt]
				
				\def\DistanceInterNiveaux{3}
				\def\DistanceInterFeuilles{0.8}
				
				\def\NiveauA{(0)*\DistanceInterNiveaux}
				\def\NiveauB{(1)*\DistanceInterNiveaux}
				\def\NiveauC{(2)*\DistanceInterNiveaux}
				\def\InterFeuilles{(-1)*\DistanceInterFeuilles}
				
				\node[noeud] (R) at ({\NiveauA},{(1.5)*\InterFeuilles}) {};
				\node[noeud] (Ra) at ({\NiveauB},{(0.5)*\InterFeuilles}) {$B_n$};
				\node[feuille] (Raa) at ({\NiveauC},{(0)*\InterFeuilles}) {$B_{n+1}$};
				\node[feuille] (Rab) at ({\NiveauC},{(1)*\InterFeuilles}) {$\overline{B_{n+1}}$};
				\node[noeud] (Rb) at ({\NiveauB},{(2.5)*\InterFeuilles}) {$\overline{B_n}$};
				\node[feuille] (Rba) at ({\NiveauC},{(2)*\InterFeuilles}) {$B_{n+1}$};
				\node[feuille] (Rbb) at ({\NiveauC},{(3)*\InterFeuilles}) {$\overline{B_{n+1}}$};
				
				\draw[fleche] (R.east)--(Ra.west) node[etiquette] {$p_n$};
				\draw[fleche] (Ra.east)--(Raa.west) node[etiquette] {$0,9$};
				\draw[fleche] (Ra.east)--(Rab.west) node[etiquette] {$0,1$};
				\draw[fleche] (R.east)--(Rb.west) node[etiquette] {$1 - p_n$};
				\draw[fleche] (Rb.east)--(Rba.west) node[etiquette] {$0,4$};
				\draw[fleche] (Rb.east)--(Rbb.west) node[etiquette] {$0,6$};
			\end{tikzpicture}
	\hfill~
		

	\item Là encore, les évènements $B_n$ et $\overline{B_n}$ partitionnent l'univers, la loi des probabilités totales donne :

$p_{n+1} = P(B_{n+1}) = P\Big(B_n \cap B_{n+1}\Big) + P\Big(\overline{B_n} \cap B_{n+1} \Big) = p_n \times 0,9 + (1-p_n) \times 0,4 = 0,9p_n + 0,4 - 0,4p_n $

Donc on a bien $p_{n+1} = 0,5p_n + 0,4$ : on a donc bien établi la relation de récurrence annoncée.

\item 
	\begin{enumerate}
		\item \emph{Initialisation :} à l'indice $n=0$, on a $p_0 = 1$ donc l'inégalité $p_0 \geqslant 0,8$ est vraie pour $n = 0$.

\emph{Hérédité :} supposons que, pour un entier naturel $n$, l'inégalité $p_n \geqslant 0,8$ est vraie.

Par hypothèse de récurrence : $\begin{aligned}[t] p_n \geqslant 0,8 &\implies 0,5p_n \geqslant 0,4 \quad \text{ car } 0,5 > 0\\
				&\implies 0,5p_n + 0,4 \geqslant 0,4 + 0,4\\
				&\implies p_{n+1} \geqslant 0,8\\ \end{aligned}$

Si l'inégalité est vérifiée à l'indice $n$, alors, elle l'est aussi au rang suivant.


\emph{Conclusion :} l'inégalité est vérifiée à l'indice 0 et sa véracité est héréditaire pour tout indice $n$ naturel, donc, par principe de récurrence, on a :

$\forall n \in \N, \quad p_n \geqslant 0,8$.

Autrement dit, on vient d'établir que la suite est minorée par 0,8.

		\item En assimilant les probabilités à des proportions, l'entreprise peut communiquer en annonçant qu'au moins 80\,\% de son parc de trottinettes est toujours en bon état.
\end{enumerate}
\item 
		\begin{enumerate}
			\item Établissons la relation de récurrence de la suite $\left(u_n\right)$. Soit $n$ un entier naturel :

$\begin{aligned} u_{n+1} &= p_{n+1} - 0,8 \quad \text{par définition de la suite } (u_n)\\
&= (0,5p_n + 0,4) - 0,8 \quad \text{par la relation de récurrence de la suite } (p_n)\\
&= 0,5p_n - 0,4\\
&= 0,5(p_n - 0,8)\\
&= 0,5u_n \quad \text{par définition de la suite } \left(u_n\right)\\ \end{aligned}$

La relation de récurrence de la suite $\left(u_n\right)$ est donc bien celle d'une suite géométrique, de raison $q = 0,5$. 

Le premier terme de la suite est $u_0 = p_0 - 0,8 = 1 - 0,8 = 0,2$.

		\item On peut donc donner la forme explicite du terme général de la suite géométrique :

$\forall n \in \N, \quad u_n = u_0 \times q^n = 0,2\times 0,5^n$.

On en déduit : $u_n = p_n - 0,8 \iff p_n = u_n + 0,8 = 0,2\times 0,5^n + 0,8$.

		\item La raison de la suite géométrique $u$ est comprise entre $- 1$ et 1, strictement, donc la suite $\left(u_n\right)$ converge vers 0.

Par limite de la somme, on en déduit que la suite $p$ converge vers $0,8$.
	\end{enumerate}
\end{enumerate}

\bigskip
\textbf{Partie B}
\medskip

	\begin{enumerate}
		\item 	Ici, nous avons :
		\begin{itemize}
			\item une expérience à deux issues (choisir une trottinette dans le parc de l'entreprise), le succès est \og choisir une trottinette en bon état \fg{}, de probabilité $p = 0,8$;
			\item cette expérience est répétée quinze fois (on constitue un lot de 15 trottinettes) d'une façon assimilable à une répétition identique et indépendante (le prélèvement des 15 trottinettes est assimilable à un tirage avec remise), donc $n=15$ répétitions;
			\item dans ce schéma de Bernoulli de paramètres $n=15$ et $p=0,8$, on note $X$ la variable aléatoire qui compte le nombre de succès (le nombre de trottinettes en bon état dans le lot).
		\end{itemize}
Avec ces éléments, on peut dire que $X$ suit la loi binomiale $ \mathcal{B}\left(15 ~;~ 0,8\right) $.

	\item La probabilité demandée est $P(X = 15)$. Par propriété, on a :

$\displaystyle P(X=15) = \binom{15}{15} \times 0,8^{15} \times 0,2^0 = 0,8^{15}$.

La probabilité que les quinze trottinettes soient en bon état est donc de $0,8^{15}$ (ici, le sujet ne précise pas de consigne d'arrondi, donc on donne la valeur exacte. La valeur approchée au dix-millième près est 0,0352).

	\item La probabilité demandée est $P(X \geqslant 10)$.

Selon le modèle de calculatrice, on peut obtenir une valeur approchée de ce résultat directement, ou alors on a recours à : $P(X \geqslant 10) = 1 - P(X <10) = 1 - P(X \leqslant 9)$.

La calculatrice donne une valeur approchée au dix-millième près qui est 0,9389.

	\item Pour une variable aléatoire suivant la loi binomiale, on a $E(X) = np = 15 \times 0,8 = 12$ (cette justification n'était pas attendue, ici).

Cela s'interprète en disant qu'en moyenne, sur un lot de quinze trottinettes choisies dans le parc de cette entreprise, douze d'entre elles seront en bon état.
\end{enumerate}

\bigskip

\end{document}
